\documentclass[CJK,13pt]{beamer}
\usepackage{CJKutf8}
\usepackage{beamerthemesplit}
\usetheme{Malmoe}
\useoutertheme[footline=authortitle]{miniframes}
\usepackage{amsmath}
\usepackage{amssymb}
\usepackage{graphicx}
\usepackage{eufrak}
\usepackage{color}
\usepackage{slashed}
\usepackage{simplewick}
\usepackage{tikz}
\usepackage{tcolorbox}
\usepackage{ulem}
\graphicspath{{../figures/}}
%%figures
\def\lfig#1#2{\includegraphics[width=#1 in]{#2}}
\def\tfig#1#2{\includegraphics[height=#1 in]{#2}}
\def\addfig#1#2{\begin{center}\includegraphics[width=#1 in]{#2}\end{center}}
\def\question{\includegraphics[width=0.3in]{why.jpg}\,}
\def\answer{\includegraphics[width=0.3in]{baozou_haha.png}\,}
\def\wulian{\includegraphics[width=0.18in]{emoji_wulian.jpg}}
\def\bigwulian{\includegraphics[width=0.35in]{emoji_wulian.jpg}}
\def\bye{\includegraphics[width=0.18in]{emoji_bye.jpg}}
\def\bigbye{\includegraphics[width=0.35in]{emoji_bye.jpg}}
\def\huaixiao{\includegraphics[width=0.18in]{emoji_huaixiao.jpg}}
\def\bighuaixiao{\includegraphics[width=0.35in]{emoji_huaixiao.jpg}}
\def\jianxiao{\includegraphics[width=0.18in]{emoji_jianxiao.jpg}}
\def\bigjianxiao{\includegraphics[width=0.35in]{emoji_jianxiao.jpg}}
\def\haoqi{\includegraphics[width=0.18in]{emoji_haoqi.jpg}}
%% colors
\def\blacktext#1{{\color{black}#1}}
\def\bluetext#1{{\color{blue}#1}}
\def\redtext#1{{\color{red}#1}}
\def\darkbluetext#1{{\color[rgb]{0,0.2,0.6}#1}}
\def\skybluetext#1{{\color[rgb]{0.2,0.7,1.}#1}}
\def\cyantext#1{{\color[rgb]{0.,0.5,0.5}#1}}
\def\greentext#1{{\color[rgb]{0,0.7,0.1}#1}}
\def\darkgray{\color[rgb]{0.2,0.2,0.2}}
\def\lightgray{\color[rgb]{0.6,0.6,0.6}}
\def\gray{\color[rgb]{0.4,0.4,0.4}}
\def\blue{\color{blue}}
\def\red{\color{red}}
\def\orange{\color[rgb]{1.,0.8,0.}}
\def\green{\color{green}}
\def\darkgreen{\color[rgb]{0,0.4,0.1}}
\def\darkblue{\color[rgb]{0,0.2,0.6}}
\def\skyblue{\color[rgb]{0.2,0.7,1.}}
%%control
\def\heaviside{\,\mathrm{h}\,}
\def\bral{\left(\begin{array}{l}}
\def\brar{\end{array}\right)}
\def\brall{\left(\begin{array}{ll}}
\def\brarr{\end{array}\right)}
\def\bralll{\left(\begin{array}{lll}}
\def\brarrr{\end{array}\right)}
\def\branchl{\left\{\begin{array}{l}}
\def\branchr{\end{array}\right.}
\def\branchll{\left\{\begin{array}{ll}}
\def\branchrr{\end{array}\right.}
\def\branchlll{\left\{\begin{array}{lll}}
\def\branchrrr{\end{array}\right.}
\def\sfgamma#1{\,\Gamma\left( #1 \right)\,}
\def\be{\begin{equation}}
\def\ee{\nonumber\end{equation}}
\def\bea{\begin{eqnarray}}
\def\eea{\nonumber\end{eqnarray}}
\def\bch{\begin{CJK}{UTF8}{gbsn}}
\def\ech{\end{CJK}}
\def\bitem{\begin{itemize}}
\def\eitem{\end{itemize}}
\def\bcenter{\begin{center}}
\def\ecenter{\end{center}}
\def\bex{\begin{minipage}{0.2\textwidth}\includegraphics[width=0.6in]{jugelizi.png}\end{minipage}\begin{minipage}{0.76\textwidth}}
\def\eex{\end{minipage}}
\def\chtitle#1{\frametitle{\bch#1\ech}}
\def\bmat#1{\left(\begin{array}{#1}}
\def\emat{\end{array}\right)}
\def\bcase#1{\left\{\begin{array}{#1}}
\def\ecase{\end{array}\right.}
\def\bmini#1{\begin{minipage}{#1\textwidth}}
\def\emini{\end{minipage}}
\def\tbox#1{\begin{tcolorbox}#1\end{tcolorbox}}
\def\pfrac#1#2#3{\left(\frac{\partial #1}{\partial #2}\right)_{#3}}
\def\res#1#2{\,\mathrm{res}\,#1\left(#2\right)\,}
\def\newt#1#2{\left(\begin{array}{c}#1\\ #2\end{array}\right)}
\def\reof#1{\,\mathrm{Re}{\left(#1\right)\,}}
\def\imof#1{\,\mathrm{Im}{\left(#1\right)\,}}
\def\Arg#1{\,\mathrm{Arg}\,#1\,}
%%symbols
\def\intfull{\int_{-\infty}^\infty}
\def\inthalf{\int_0^\infty}
\def\bropt{\,(\ \ \ )}
\def\sech{\mathrm{sech}\,}
\def\csch{\mathrm{csch}\,}
\def\asinh{\mathrm{asinh}\,}
\def\acosh{\mathrm{acosh}\,}
\def\sone{$\star$}
\def\stwo{$\star\star$}
\def\sthree{$\star\star\star$}
\def\sfour{$\star\star\star\star$}
\def\sfive{$\star\star\star\star\star$}
\def\rint{{\int_\leftrightarrow}}
\def\roint{{\oint_\leftrightarrow}}
\def\stdHf{{\textit{\r H}_f}}
\def\deltaH{{\Delta \textit{\r H}}}
\def\ii{{\dot{\imath}}}
\def\skipline{{\vskip0.1in}}
\def\skiplines{{\vskip0.2in}}
\def\lagr{{\mathcal{L}}}
\def\hamil{{\mathcal{H}}}
\def\vecT{{\mathbf{T}}}
\def\vecN{{\mathbf{N}}}
\def\vecB{{\mathbf{B}}}
\def\vecv{{\mathbf{v}}}
\def\vecr{{\mathbf{r}}}
\def\vecf{{\mathbf{f}}}
\def\vecg{{\mathbf{g}}}
\def\vecupsilon{{\mathbf{\upsilon}}}
\def\vecu{{\mathbf{u}}}
\def\vecj{{\mathbf{j}}}
\def\vecx{{\mathbf{x}}}
\def\vecy{{\mathbf{y}}}
\def\vecz{{\mathbf{z}}}
\def\veck{{\mathbf{k}}}
\def\vecp{{\mathbf{p}}}
\def\vecn{{\mathbf{n}}}
\def\vecA{{\mathbf{A}}}
\def\vecP{{\mathbf{P}}}
\def\vecsigma{{\mathbf{\sigma}}}
\def\hatJn{{\hat{J_\vecn}}}
\def\hatJx{{\hat{J_x}}}
\def\hatJy{{\hat{J_y}}}
\def\hatJz{{\hat{J_z}}}
\def\hatj#1{\hat{J_{#1}}}
\def\hatphi{{\hat{\phi}}}
\def\hatq{{\hat{q}}}
\def\hatpi{{\hat{\pi}}}
\def\vel{\upsilon}
\def\Dint{{\mathcal{D}}}
\def\adag{{\hat{a}^\dagger}}
\def\bdag{{\hat{b}^\dagger}}
\def\cdag{{\hat{c}^\dagger}}
\def\ddag{{\hat{d}^\dagger}}
\def\hata{{\hat{a}}}
\def\hatb{{\hat{b}}}
\def\hatc{{\hat{c}}}
\def\hatd{{\hat{d}}}
\def\hatD{{\,\hat{D}}}
\def\hatN{{\hat{N}}}
\def\hatH{{\hat{H}}}
\def\hatp{{\hat{p}}}
\def\Fup{{F^{\mu\nu}}}
\def\Fdown{{F_{\mu\nu}}}
\def\newl{\nonumber \\}
\def\vece{\mathrm{e}}
\def\calM{{\mathcal{M}}}
\def\calT{{\mathcal{T}}}
\def\calR{{\mathcal{R}}}
\def\barpsi{\bar{\psi}}
\def\baru{\bar{u}}
\def\barv{\bar{\upsilon}}
\def\qeq{\stackrel{?}{=}}
\def\ftf{\stackrel{\mathcal{FT}}{\Longrightarrow}}
\def\ftb{\stackrel{\mathcal{FT}}{\Longleftarrow}}
\def\ftfb{\stackrel{\mathcal{FT}}{\Longleftrightarrow}}
\def\ltf{\stackrel{\mathcal{LT}}{\Longrightarrow}}
\def\ltb{\stackrel{\mathcal{LT}}{\Longleftarrow}}
\def\ltfb{\stackrel{\mathcal{LT}}{\Longleftrightarrow}}
\def\torder#1{\mathcal{T}\left(#1\right)}
\def\rorder#1{\mathcal{R}\left(#1\right)}
\def\contr#1#2{\contraction{}{#1}{}{#2}#1#2}
\def\trof#1{\mathrm{Tr}\left(#1\right)}
\def\trace{\mathrm{Tr}}
\def\comm#1{\ \ \ \left(\mathrm{used}\ #1\right)}
\def\tcomm#1{\ \ \ (\text{#1})}
\def\slp{\slashed{p}}
\def\slk{\slashed{k}}
\def\calp{{\mathfrak{p}}}
\def\veccalp{\mathbf{\mathfrak{p}}}
\def\Tthree{T_{\tiny \textcircled{3}}}
\def\pthree{p_{\tiny \textcircled{3}}}
\def\dbar{{\,\mathchar'26\mkern-12mu d}}
\def\erf{\mathrm{erf}}
\def\const{\mathrm{const.}}
\def\pheat{\pfrac p{\ln T}V}
\def\vheat{\pfrac V{\ln T}p}
%%units
\def\fdeg{{^\circ \mathrm{F}}}
\def\cdeg{^\circ \mathrm{C}}
\def\atm{\,\mathrm{atm}}
\def\angstrom{\,\text{\AA}}
\def\SIL{\,\mathrm{L}}
\def\SIkm{\,\mathrm{km}}
\def\SIyr{\,\mathrm{yr}}
\def\SIGyr{\,\mathrm{Gyr}}
\def\SIV{\,\mathrm{V}}
\def\SImV{\,\mathrm{mV}}
\def\SIeV{\,\mathrm{eV}}
\def\SIkeV{\,\mathrm{keV}}
\def\SIMeV{\,\mathrm{MeV}}
\def\SIGeV{\,\mathrm{GeV}}
\def\SIcal{\,\mathrm{cal}}
\def\SIkcal{\,\mathrm{kcal}}
\def\SImol{\,\mathrm{mol}}
\def\SIN{\,\mathrm{N}}
\def\SIHz{\,\mathrm{Hz}}
\def\SIm{\,\mathrm{m}}
\def\SIcm{\,\mathrm{cm}}
\def\SIfm{\,\mathrm{fm}}
\def\SImm{\,\mathrm{mm}}
\def\SInm{\,\mathrm{nm}}
\def\SImum{\,\mathrm{\mu m}}
\def\SIJ{\,\mathrm{J}}
\def\SIW{\,\mathrm{W}}
\def\SIkJ{\,\mathrm{kJ}}
\def\SIs{\,\mathrm{s}}
\def\SIkg{\,\mathrm{kg}}
\def\SIg{\,\mathrm{g}}
\def\SIK{\,\mathrm{K}}
\def\SImmHg{\,\mathrm{mmHg}}
\def\SIPa{\,\mathrm{Pa}}
%page
\def\secpage#1#2{\begin{frame}\bcenter{\bf \Huge #1} \ecenter \skipline \tbox{\bcenter #2 \ecenter}\end{frame}}
\def\append#1#2{\secpage{附录#1}{#2}}
\def\thinka#1{\begin{frame}\frametitle{思考题}\bcenter\lfig{0.4}{think0.jpg}  \skipline #1 \ecenter\end{frame}}
\def\thinkb#1{\begin{frame}\frametitle{思考题}\bcenter\lfig{0.5}{think1.jpg}  \skipline #1 \ecenter\end{frame}}
\def\thinkc#1{\begin{frame}\frametitle{思考题}\bcenter\lfig{0.5}{think2.jpg}  \skipline #1 \ecenter\end{frame}}
\def\thinkd#1{\begin{frame}\frametitle{思考题}\bcenter\lfig{0.5}{think3.jpg}  \skipline #1 \ecenter\end{frame}}
\def\thinke#1{\begin{frame}\frametitle{思考题}\bcenter\lfig{0.5}{think4.jpg}  \skipline #1 \ecenter\end{frame}}
\def\thinkf#1{\begin{frame}\frametitle{思考题}\bcenter\lfig{0.5}{think5.jpg}  \skipline #1 \ecenter\end{frame}}

\def\courseurl{http://zhiqihuang.top/gr}

\def\tpage#1#2{
\title{GR \S{#1}  #2}
  \author{Zhiqi Huang}

\begin{frame}
\begin{center}
{\bf \Huge G}eneral {\bf \Huge R}elativity

{\vskip 0.1in}



{\Large \S #1 #2}

{\vskip 0.2in}

{Lecturer: 黄志琦}

\vskip 0.2in

\courseurl

\end{center}
\end{frame}
}


  \date{}
  \begin{document}
  \bch
\tpage{7}{Covariant Derivative}


\begin{frame}
  在这一讲里我们继续考虑三维欧氏空间里的曲面 $\vecx(u^1, u^2)$ ——虽然很多讨论并不需要在如此特殊的假设下进行，但

  \skiplines

  \tbox{这门课的目的是尽可能地让憨憨们理解广义相对论。}
\end{frame}


\begin{frame}
  \frametitle{曲面上的标量场}
  \addfig{2}{Efield.jpg}
  
  我们想象三维欧氏空间中有静电势 $\varphi(\vecx)$。在曲面上的爬虫也能够测量曲面上的静电势 $\varphi(u^1, u^2)$。

\end{frame}

\begin{frame}
\frametitle{标量场的坐标偏导是协变矢量}

电势的偏导数 $\varphi_{,i}$ 在坐标变换 $(u^1,u^2)\rightarrow (\tilde{u}^1, \tilde{u}^2)$ 下满足张量的变换规则：
$$ \widetilde{\varphi_{,i}} = \frac{\partial u^j}{\partial \tilde{u}^i} \varphi_{,j}.$$
这说明标量场的坐标偏导数是矢量。它的指标在下面，我们称之为{\blue 协变矢量(covariant vector)}。指标在上面的矢量叫{\blue 逆变矢量(contravariant vector)}。

\bmini{0.33}
\lfig{1.2}{qiubite.jpg}
\emini
\bmini{0.63}
注意：{\blue 矢量是客观的物理存在，协变形式和逆变形式只是它们在给定坐标系里的不同形式的投影。}

{\scriptsize 啥，你问为什么要用“协变矢量”这样有误导性的名词？“矢量的协变形式”是7个字，“协变矢量”是4个字，我懒得跟你说下去了……}
\emini
\end{frame}


\begin{frame}

  
  \frametitle{我们来讨论矢量的协变形式和逆变形式到底是怎么投影的——}

  \addfig{1}{biehuang.jpg}
  
  别慌，我没打算给你讲切丛，余切丛，微分形式、抽象指标、旋量表示……
  
\end{frame}


\begin{frame}
  \frametitle{简单粗暴一句话可以解释}
  \tbox{逆变分解，协变投影。}

  \skipline
  
  好的我知道你木有看懂这句话，我们得上图——
\end{frame}


\begin{frame}
  \frametitle{当 $u^1, u^2$ 是弧长元变量}
  \bmini{0.42}
  \lfig{1.8}{covcontra.jpg}

  \skipline
  
  {\small 我们在学习曲线论时，发现使用弧长变量可以大大简化几何图像。在曲面论里也是如此：  }
  \emini
  \hspace{0.1in}
  \bmini{0.5}
  {\small 假设 $(u^1, u^2)$ 在 $O$ 点是弧长变量（即 $O$ 点的 $g_{11}= g_{22}= 1$），图中标出了变化坐标 $u^1, u^2$ 产生的两个单位切矢的方向。}

  红色的箭头是在 $O$ 点的一个矢量 $\vec{A}$ （物理客观存在）。
  
  \skipline
  
{\blue  逆变形式 $(A^1, A^2)$ 是 $\vec{A}$ 按两个单位切矢分解的系数——或者说，它们是按平行四边形法则投影的。}

  \skipline

 {\blue 协变形式 $(A_1, A_2)$ 是 $\vec{A}$ 在两个单位切矢上的投影——对，就是简单粗暴的作垂线投影！}
  
  
  \emini
\end{frame}


\begin{frame}
  \frametitle{一般情况要做下长度单位的修正}
  \addfig{2}{covcontra2.jpg}
\end{frame}

\begin{frame}
  \frametitle{代数表达式}
  当你明白了协变和逆变矢量都是依赖于坐标系的特定投影方式之后，前面唠叨的那堆就不那么重要了，具体计算时可以使用三维欧氏空间的图像来辅助：

  \tbox{  $$ \vec{A} = A^i\vecx_{,i} = A_i\vecx^{,i}$$
    $$ A_i = \vec{A} \cdot \vecx_{,i} ,   A^i = \vec{A}\cdot\vecx^{,i}$$}

  这些公式可以用我们之前用过的秘密武器轻松证明。证明时注意 $\vec{A}$ 被限定在曲面上，没有法向分量。

  \bmini{0.3}
  \addfig{1}{zhongdian2.jpg}
  \emini
  \bmini{0.65}
        {\blue 我们说曲面上的客观物理对象（张量）都是指限制在曲面内的物理量（例如矢量必须沿曲面的切向）。}
        \emini
\end{frame}



\begin{frame}
  
\addfig{0.8}{eg.jpg}
  
因为 $\varphi_{,i}$ 是静电势在曲面上的梯度（客观物理对象）在第 $i$ 个坐标方向上的投影——重点是，它是投影——所以 $\varphi_{,i}$ 是协变矢量。


\skiplines

{\scriptsize 说因为  $\varphi_{,i}$ 指标在下面所以它是协变矢量的憨憨——你说的都对，但我觉得你应该考虑一下退课……}

\end{frame}


\begin{frame}
  
\addfig{0.8}{ega.jpg}
  
曲面上位移矢量 $d\vecx$ 按照两个自然切向量的分解式为
$$d\vecx = \vecx_{,i}du^i$$
重点是， 它是分解系数。所以 $du^i$ 是逆变矢量。

\skiplines

{\scriptsize 说因为  $du^i$ 指标在上面所以它是逆变矢量的憨憨——你说的都对，而且我刚刚想起来你好像来不及退课了……}

\end{frame}

\begin{frame}
  \frametitle{矢量场的微分}
  搞清楚协变和逆变的几何图像之后，我们来研究更复杂的对象： 曲面上的矢量场 $\vec{A}$ 的微分。这里符号 $\vec{A}$ 是指物理客观存在。

  比如，对很近的两点 $P$ 和 $Q$，我们希望计算物理量
  $$ d\vec{A} = \vec{A}(P) - \vec{A}(Q) $$
  注意它既不能用
  $$ A_i(P) - A_i(Q).$$
  也不能用
  $$ A^i(P) - A^i(Q).$$  
  来描述！

  因为对随意取的曲面坐标系， $P$ 处的投影规则和 $Q$ 处的投影规则可能不一样，这样求差就相当于把投影规则的差别也囊括进来了。  
  
\end{frame}


\begin{frame}
  \frametitle{矢量场的微分（续）}
  为了能用数字描述物理量，我们还是要取一个投影规则。但绝不是两个！

  \skipline

  如果都用在 $Q$ 处的投影规则，
  
 {\blue $d\vec{A}$ 的协变形式 = $\vec{A}(P)$ 在 $Q$ 的协变投影 - $\vec{A}(Q)$ 的协变形式}
 
 如果你都用 $P$ 点的投影规则，
 
 {\blue $d\vec{A}$ 的协变形式 = $\vec{A}(P)$ 的协变形式 - $\vec{A}(Q)$ 在 $P$ 的协变投影} 

 这两种定义只有一个二阶小量的差别，并不会产生什么问题。重点是你要坚持只用一个投影规则，不要把投影规则的差异（一阶小量）错误地加到 $d\vec{A}$ 里来。
  
  \skipline
  
  当然对 $d\vec{A}$ 的逆变形式，类似的讨论也都成立。
\end{frame}


\begin{frame}
  \frametitle{矢量场的微分（续）}
   借助三维欧氏空间的辅助，我们可以算出 $d\vec{A}$ 的协变形式为
  \bea
  (dA)_i &=& (d \vec{A}) \cdot \vecx_{,i} \newl
  &=& d(\vec{A}\cdot\vecx_{,i}) - \vec{A}\cdot d\vecx_{,i} \newl
  &=& d(A_i) - \left(A_k\vecx^{,k}\right)\cdot\left(\vecx_{,i,j}du^j\right) \newl
  &=& d(A_i) - \Gamma^k_{\ ij} A_k du^j     
  \eea
\end{frame}


\begin{frame}
  \frametitle{矢量场的协变微商}
  有了微分 $d\vec{A}$ 的正确表述形式，就可以得到矢量场的正确的求导姿势——称作{\blue 协变微商(covariant derivative)}。
  {\blue $$ A_{i;j} = A_{i,j} - \Gamma^k_{\ ij} A_k .$$
  为了区别用逗号(comma)表示的普通偏导数，我们用分号(semicolon)来表示协变微商。}
  
  \addfig{1}{zhongdian2.jpg}
  
  “协变”是指求完导之后多了个下指标，和被求导的张量是用什么形式表述的没有关系；“微商”和朋友圈没关系。
\end{frame}



\begin{frame}
  \frametitle{协变微商满足普通微商的所有性质}

  从前面的讨论看出来，{\blue 协变微商才是“正确的微商”，它满足微商的所有性质，并且张量的协变微商还是张量}。普通微商的物理意义比较复杂，好在它们的数学形式比较简单，我们可以撇开其物理意义，直接用数学知识进行操作。要注意的是，张量的普通微商一般不是张量（不涉及投影的标量除外）。

  \addfig{1}{hairgone.jpg}
  
  这是一个你同时精通数学和物理\sout{才能保住}也保不住头发的绝佳例子
\end{frame}


\begin{frame}
  \frametitle{逆变矢量的微商}
  逆变矢量的微商直接怼起来稍有些繁琐，但可以用下面的办法巧妙地赖出来：

  对逆变 $A^i$ 以及任意协变 $B_i$，两者可以做张量积再收缩得到一个标量 $A^iB_i$。因为标量不需要投影，所以普通微商和协变微商是一样的。即有
  $$(A^iB_i)_{;j} = (A^iB_i)_{,j} = A^i_{\ ,j}B_i + A^i B_{i,j}.$$
  另一方面
  \be
  (A^iB_i)_{;j}=  A^i_{\ ;j}B_i + A^i B_{i;j} = A^i_{\ ;j}B_i + A^iB_{i,j} - A^i \Gamma^k_{\ ij}B_k
  \ee
  比较两个结果，即有
  \be
  A^i_{\ ;j}B_i = A^i_{\ ,j}B_i +  A^i \Gamma^k_{\ ij}B_k
  \ee
\end{frame}

\begin{frame}
  \frametitle{逆变矢量的微商(续）}
    由于求和指标可以随便替换，上述结果中最后一项把 $i, k$ 互换，得到
  \be
  A^i_{\ ;j}B_i = A^i_{\ ,j}B_i +  A^k \Gamma^i_{\ kj}B_i
  \ee
  上式要对任意 $B_i$ 成立，所以{\blue
  \be
  A^i_{\ ;j} = A^i_{\ ,j} +  \Gamma^i_{\ kj}A^k 
  \ee
  }
\end{frame}

\begin{frame}
  \frametitle{高阶张量的协变微商}
  一般来说，大多数高阶张量都不能写成矢量的张量积。但是它们都能拆成一堆“矢量的张量积”之和。例如对一般的二阶张量 $T_{ij}$，我们不能指望存在 $A_i, B_j$ 使得 $T_{ij}  =A_iB_j$，但是我们可以假设
  $$ T_{ij} = A_iB_j + C_i D_j + \ldots$$
  这个结论无非是在说“给了足够多自由度方程总有解”，理解起来并不困难。

  \skipline
  
  每一项 $A_iB_j$ 形式的二阶张量的协变微商可以用矢量的协变微商法则推导出来。最后叠加的效果是 

 {\blue $$ T_{ij;k} = T_{ij, k} - \Gamma^l_{ik} T_{lj} - \Gamma^l_{jk} T_{il}.$$}

\end{frame}


\begin{frame}
  \frametitle{高阶张量的协变微商（续）}
  对一般的带一堆指标的张量形式，协变微商的计算规则如下：
  {\blue
  \bea
  \left(T^{i_1i_2\ldots}_{\ \ j_1j_2\ldots}\right)_{;k} &=& \left(T^{i_1i_2\ldots}_{\ \ j_1j_2\ldots}\right)_{,k} \newl
  && + \Gamma^{i_1}_{\ lk} T^{li_2\ldots}_{\ \ j_1j_2\ldots} + \Gamma^{i_2}_{\ lk} T^{i_1l\ldots}_{\ \ j_1j_2\ldots} + \ldots \newl
  && - \Gamma^l_{j_1k}T^{i_1i_2\ldots}_{\ \ lj_2\ldots} - \Gamma^l_{j_2k}T^{i_1i_2\ldots}_{\ \ j_1l\ldots}-\ldots
  \eea
  }
\end{frame}


\begin{frame}
  \frametitle{度规只会跟你聊规则}

  度规是一个很特殊的二阶张量，它描述了协变和逆变的投影规则。

  \addfig{1.5}{tangmetric.jpg}

  除此之外，它啥也不干。
  
\end{frame}


\begin{frame}
  \frametitle{度规的协变微商是零}
  
  一般来说，度规 $g_{ij}$ 在各处并不相同。但这完全是由于各处投影规则的变化产生的。由于协变微商里已经剔除了投影规则变化的影响，所以我们推测——度规的协变微商应该处处是零。

  \addfig{0.8}{caozuo.jpg}
  
  证明如下: 因为 $g^i_{\ j}$ 是单位矩阵，普通微商消失，所以
  $$ g^i_{\ j;k} = \Gamma^i_{lk}g^l_{\ j} - \Gamma^l_{jk}g^i_{\ l} = \Gamma^i_{\ jk} - \Gamma^i_{\ jk} = 0. $$
  又协变微商有保张量性，所以 $g^i_{\ j;k}$ 是张量，可以把指标 $i$ 降下来 或者把 指标 $j$ 升上去。所以无论度规以何种形式写出来，它的协变微商总是零。
  
\end{frame}


\begin{frame}
  \frametitle{对度规的协变微商是零的理解}
  物理上讲，度规描述的是度量规则。
  \bitem
\item{对嵌入高维欧氏空间的曲面，度量规则其实各处都一样（都是同一个欧氏空间度量规则），所以度规的协变微商是零。}
\item{如果你想追随高斯的脑洞，在曲面论中彻底摒弃高维空间，你就必须得承认生活在曲面上的爬虫有一把依靠物理学获得的，不依赖于曲面性质的绝对度量尺子（例如普朗克长度）来定义各处的度量规则。在此前提下，当然各处的度量规则在物理上还是一样的。}
  \eitem

  {\scriptsize 对纯粹的数学家而言，实验可检验性毫不重要。建立内蕴几何需要物理学这件事对数学家而言显然非常荒谬。当你问：爬虫到底该选取什么尺子来进行长度测量？数学家会觉得你该吃药了……}
\end{frame}




\ech
\end{document}
